
\section{External inputs}\label{app:external}

\paragraph{What this appendix proves and how to read it.}
The proof relies on a small number of published results that it does not
prove. Three standard ones are stated in \cref{sec:notation}: Weil's bound
for complete character sums (\cref{thm:weil}), used throughout, and the
explicit Chebyshev estimate for primes in the residue classes modulo four
(\cref{thm:bmor}) and Kassabov's inequality for angles between intersections
of subspaces (\cref{thm:kassabov}), whose uses are named there;
\cref{subsec:external:sources} matches them with their sources. The
second-order inputs are stated here, and each is used only where named. The
character table of $\mathrm{GL}_2$ over a finite field (\cref{thm:prelim:gl2})
enters only \cref{lem:dl:group}, hence the four-cycle kernel bound
\cref{thm:dl:K}. The $\ell$-adic inputs
(\cref{subsec:external:sheaves}) enter \cref{lem:dl:sheaf,lem:tame-trace-bound} and, through them, the
proofs of the two-variable character sum bound \cref{thm:dl:H2} and of the
moving-root moments \cref{lem:moving-root-moments}; parts (v) and (vi) of
\cref{thm:prelim:ks} are used only in the last, through
\cref{lem:dl:sheaf}(v),(vi). Through \cref{thm:dl:K,thm:dl:H2} and the
moving-root moments the second-order inputs reach the spectral caps of the
middle and upper ranges (among them, through \cref{cor:raw-pair-maps}, the caps
of the positive roots), the direct-star root
(\cref{lem:record:direct-star-inputs}), the dense comparisons and the six-query
mixture of the upper range (through \cref{lem:wt:stars}) and the conditional
rows of the upper range (\cref{subsec:advance:product-rows}). The modular input
(\cref{thm:ahlgren-ono}) enters only the second moment of the neighbourhood
count (\cref{lem:record:weighted-moment}), and through it the root-half
variances $\chi_p$ of the middle roots (in the restorations and in the
completion of the root cut L20), the star mean of that cut
(\cref{lem:l20-star-mean}), and the linear prices and the normalization of the
positive roots (\cref{lem:root-linear,lem:root-normalization});
\cref{tab:second-order-uses} lists every use. The character table is a finite
statement, and the modular input is an identity at each prime together with
Deligne's bound for the coefficients of one modular form; the $\ell$-adic inputs
concern families of algebraic curves. The proofs using them are separate from
the finite arithmetic certificates.

Peyre's tensorization theorem is not an external input: \cref{app:peyre}
proves the finite case of \cite[Theorem~3.3.1]{Peyre} that the paper uses
(\cref{thm:peyre-finite}). Together with the three theorems named above, the
results of this appendix are the external theorems on which the proof
depends; everything else is proved in the paper or is standard textbook
material, such as Khintchine's inequality, the Chinese remainder theorem and
Fourier analysis on finite abelian groups.

\emph{Notation.} The results below are quoted in the notation of their
sources, which is local to this appendix and to the proofs that apply them.
In particular $q$ is the number of elements of a finite field (the
applications take $q=p$) or, in \cite{BMOR18}, a modulus, or, in
\cref{thm:ahlgren-ono}, $e^{2\pi iz}$; $N\zeta$ is a norm
and $N(x)$ the size of a residue field; $E$ is a finite field and $t$ is the
parameter of a family of curves.
Deligne's field isomorphism $\overline{\mathbb Q}_\ell\to\C$, written
$\iota$ in the sources, is written $\jmath$, since $\iota$ is the reserved cut
of the paper; the proofs that apply these results do not name it.

\subsection{The characters of \texorpdfstring{$\mathrm{GL}_2(\F_q)$}{GL2(Fq)}}
\label{subsec:external:gl2}

For a field $\F_q$ we write $\F_{q^2}$ for its quadratic extension and, for
$\zeta\in\F_{q^2}$, $\bar\zeta=\zeta^q$ and $N\zeta=\zeta\bar\zeta$.
The group $\mathrm{GL}_2(\F_q)$ consists of the invertible $2\times2$
matrices over $\F_q$. A complex representation assigns a linear operator to
each group element, respecting products; its \emph{character} is the trace
of that operator. It is \emph{irreducible} if it has no proper nonzero
invariant subspace. The \emph{degree} of a character is the dimension of its
representation. The projective line $\mathbb P^1(\F_q)$ is the set of
one-dimensional subspaces of $\F_q^2$. For general finite-group
representation theory see \cite{Serre77,Isaacs76}; the finite-field
$\mathrm{GL}_2$ table is also treated in \cite{PiatetskiShapiro83}. We cite
the precise table as in \cite[\S5.2]{FultonHarris}.

\begin{theorem}[The characters of $\mathrm{GL}_2(\F_q)$; {\cite[\S5.2]{FultonHarris}}]\label{thm:prelim:gl2}
Let $q$ be a power of an odd prime and $\F=\F_q$. Every element of
$\mathrm{GL}_2(\F)$ is of exactly one of the following types:
\textup{(a)} a scalar matrix $xI$, $x\in\F^\times$; \textup{(b)} a
non-scalar matrix with characteristic polynomial $(T-x)^2$, $x\in\F^\times$;
\textup{(c)} a matrix with two eigenvalues $x\ne y$ in $\F^\times$;
\textup{(d)} a matrix with eigenvalues $\zeta,\bar\zeta$,
$\zeta\in\F_{q^2}\setminus\F$. Two elements of the same type with the same
eigenvalues are conjugate. For characters $\alpha,\beta$ of $\F^\times$ and
$\varphi$ of $\F_{q^2}^\times$ there are complex representations $U_\alpha$
(of dimension $1$, $g\mapsto\alpha(\det g)$), $V_\alpha$ (dimension $q$),
$W_{\alpha,\beta}$ for $\alpha\ne\beta$ (dimension $q+1$) and $X_\varphi$ for
$\varphi\ne\varphi^q$ (dimension $q-1$) such that:
\begin{enumerate}[label=\textup{(\roman*)}]
\item $V_1$ is the complement of the trivial representation in the
  permutation representation of $\mathrm{GL}_2(\F)$ on $\mathbb P^1(\F)$, and
  $V_\alpha=V_1\otimes U_\alpha$;
\item these representations are irreducible, $W_{\alpha,\beta}\cong
  W_{\beta,\alpha}$ and $X_\varphi\cong X_{\varphi^q}$, there are no other
  isomorphisms among them, and every irreducible complex representation of
  $\mathrm{GL}_2(\F)$ is isomorphic to one of them;
\item their characters take the following values at the elements of the four
  types:
\begin{center}\small
\begin{tabular}{@{}lllll@{}}
\toprule
 & (a) $xI$ & (b) & (c) $x\ne y$ & (d) $\zeta,\bar\zeta$\\
\midrule
$U_\alpha$ & $\alpha(x)^2$ & $\alpha(x)^2$ & $\alpha(x)\alpha(y)$ & $\alpha(N\zeta)$\\
$V_\alpha$ & $q\,\alpha(x)^2$ & $0$ & $\alpha(x)\alpha(y)$ & $-\alpha(N\zeta)$\\
$W_{\alpha,\beta}$ & $(q+1)\alpha(x)\beta(x)$ & $\alpha(x)\beta(x)$ & $\alpha(x)\beta(y)+\alpha(y)\beta(x)$ & $0$\\
$X_\varphi$ & $(q-1)\varphi(x)$ & $-\varphi(x)$ & $0$ & $-(\varphi(\zeta)+\varphi(\bar\zeta))$\\
\bottomrule
\end{tabular}
\end{center}
\item the conjugacy class of an element of type \textup{(a)}, \textup{(b)},
  \textup{(c)}, \textup{(d)} has $1$, $q^2-1$, $q(q+1)$, $q(q-1)$ elements,
  and there are $q-1$, $q-1$, $(q-1)(q-2)/2$, $q(q-1)/2$ classes of the
  respective types. The centralizer of a non-scalar element $g$ is the
  group of units of the algebra $\F[g]$, of order $q(q-1)$, $(q-1)^2$,
  $q^2-1$ for $g$ of type \textup{(b)}, \textup{(c)}, \textup{(d)}; the
  group $\mathrm{GL}_2(\F)$ has order $(q^2-1)(q^2-q)$.
\end{enumerate}
\end{theorem}

\noindent Part (iv) is the class table at the beginning of
\cite[\S5.2]{FultonHarris}. The centralizers named there (upper triangular
matrices with equal diagonal entries, the diagonal torus, and
$\F_{q^2}^\times$ acting on $\F_{q^2}\cong\F^2$) are the unit groups of
$\F[g]$ for the listed representatives, and the class sizes are their
indices.

\subsection{\texorpdfstring{$\ell$}{l}-adic sheaves on curves}
\label{subsec:external:sheaves}

We use constructible $\overline{\mathbb Q}_\ell$-sheaves and \'etale
cohomology with compact supports on schemes of finite type over a finite
field $\F_q$ of odd characteristic $p$, as in \cite{SGA4half} and
\cite[\S1]{DeligneWeilII}. The exponential-sum and monodromy background is
developed in Katz's books \cite{KatzGKM,KatzESDE90,KatzTLM02,KatzMMP05}. Here $\ell\ne p$ is a
fixed prime, and we fix a field isomorphism
$\jmath:\overline{\mathbb Q}_\ell\to\C$, through which traces are regarded as
complex numbers. In the applications (\cref{thm:dl:H2} and the moving-root
moments) these facts are applied to dense open subschemes of the affine
line; the remark after \cref{thm:prelim:ks} also applies them to the fibres
of the Katz--Sarnak family.
A \emph{lisse} sheaf is locally constant in the \'etale topology and has
finite rank; its stalk is the vector space attached to a geometric point,
and its rank is the dimension of that space. A \emph{constructible} sheaf
admits a finite stratification on whose strata it is lisse. The cohomology
group $H^i_c$ denotes \'etale cohomology with compact supports. For a group
$G$ acting on a vector space $W$, its space of \emph{coinvariants} is
$W_G=W/\langle gw-w:g\in G,w\in W\rangle$. The Tate twist $(-1)$ is the
one-dimensional factor on which geometric Frobenius over $\F_q$ acts by $q$.
We use the following standard facts.
\begin{enumerate}[label=\textup{(E\arabic*)},leftmargin=*]
\item\label{it:prelim:E1} \emph{Lisse sheaves and monodromy.} Let $X$ be a
  smooth geometrically connected curve over $\F_q$ with a geometric generic
  point $\bar\eta$. A lisse sheaf $\mathcal F$ on $X$ is the same as a
  continuous representation $\rho_{\mathcal F}$ of $\pi_1(X,\bar\eta)$ on the
  finite-dimensional space $\mathcal F_{\bar\eta}$
  \cite[(1.1.6)]{DeligneWeilII}. For $x\in X(\F_q)$ the geometric Frobenius
  $\mathrm{Frob}_x\in\pi_1(X,\bar\eta)$ is defined up to conjugation
  \cite[(1.1.7), (1.1.8)]{DeligneWeilII}, and
  $\mathrm{Tr}(\mathrm{Frob}_x\mid\mathcal F)$ denotes the trace of
  $\rho_{\mathcal F}(\mathrm{Frob}_x)$, i.e.\ the trace of $\mathrm{Frob}_x$
  on the stalk at a geometric point over $x$. The geometric fundamental group
  $\pi_1^{\rm geom}(X)=\pi_1(X\otimes\bar\F_q,\bar\eta)$ is a closed normal
  subgroup of $\pi_1(X,\bar\eta)$ \cite[(1.1.13), (1.1.15)]{DeligneWeilII};
  $\mathcal F$ is \emph{geometrically irreducible} if
  $\rho_{\mathcal F}|_{\pi_1^{\rm geom}(X)}$ is irreducible. If
  $X'\subseteq X$ is a dense open subscheme, the map
  $\pi_1^{\rm geom}(X')\to\pi_1^{\rm geom}(X)$ is surjective:
  $X\otimes\bar\F_q$ is connected and normal, so by
  \cite[Exp.~V, Prop.~8.2]{SGA1} already the fundamental group of its generic
  point maps onto $\pi_1^{\rm geom}(X)$, and this map factors through
  $\pi_1^{\rm geom}(X')$. Hence $\mathcal F|_{X'}$ is geometrically
  irreducible if $\mathcal F$ is. If $\mathcal L$ is lisse of rank one, a
  subspace of $\mathcal F_{\bar\eta}\otimes\mathcal L_{\bar\eta}$ of the form
  $W\otimes\mathcal L_{\bar\eta}$ is stable under
  $\rho_{\mathcal F}\otimes\rho_{\mathcal L}$ iff $W$ is stable under
  $\rho_{\mathcal F}$, so $\mathcal F\otimes\mathcal L$ is geometrically
  irreducible iff $\mathcal F$ is; and
  $\mathrm{Tr}(\mathrm{Frob}_x\mid\mathcal F\otimes\mathcal L)=
  \mathrm{Tr}(\mathrm{Frob}_x\mid\mathcal F)\,
  \mathrm{Tr}(\mathrm{Frob}_x\mid\mathcal L)$. If $\mu:X'\to X$ is an
  isomorphism of $\F_q$-schemes, $\mu^*\mathcal F$ is lisse, geometrically
  irreducible iff $\mathcal F$ is, and
  $\mathrm{Tr}(\mathrm{Frob}_x\mid\mu^*\mathcal F)=
  \mathrm{Tr}(\mathrm{Frob}_{\mu(x)}\mid\mathcal F)$.
\item\label{it:prelim:E2} \emph{Weights}
  \cite[(1.2.1), (1.2.2), (1.2.5), (1.2.6)]{DeligneWeilII}. A lisse
  $\mathcal F$ on $X$ is \emph{$\jmath$-pure of weight $w$} if for every
  closed point $x$ of $X$, with residue field of $N(x)$ elements, every
  eigenvalue $\alpha$ of $\mathrm{Frob}_x$ on $\mathcal F_{\bar x}$
  satisfies $|\jmath(\alpha)|=N(x)^{w/2}$, and \emph{pure of weight $w$} if it
  is $\jmath$-pure of weight $w$ for every choice of $\jmath$. A sheaf is
  \emph{$\jmath$-mixed of weights $\le w$} if it has a finite filtration
  whose successive quotients are $\jmath$-pure of weights $\le w$; a pure
  sheaf of weight $w$ is $\jmath$-mixed of weights $\le w$. The tensor
  product of lisse sheaves pure of weights $w$, $w'$ is pure of weight
  $w+w'$, and $\mu^*$ ($\mu$ an isomorphism) preserves purity.
\item\label{it:prelim:E3} \emph{Tame ramification} \cite[Ch.~1]{KatzGKM}.
  Let $X\subseteq\mathbb P^1$ be a dense open subscheme over $\F_q$. For a
  point $x\in\mathbb P^1(\bar\F_q)\setminus X(\bar\F_q)$ there is an
  \emph{inertia group} $I_x\subseteq\pi_1^{\rm geom}(X)$ at $x$, a closed
  subgroup well defined up to conjugation, with a unique pro-$p$ Sylow
  subgroup $P_x$ (the \emph{wild inertia group}), which is normal in $I_x$.
  The restriction of $\rho_{\mathcal F}$ to $I_x$ is the \emph{local
  monodromy} of a lisse $\mathcal F$ at $x$. A lisse $\mathcal F$ on $X$ is
  \emph{tame at $x$} if $\rho_{\mathcal F}(P_x)=1$, and then its Swan
  conductor $\mathrm{Swan}_x(\mathcal F)$ is $0$; its local monodromy then
  factors through the tame quotient $I_x/P_x$, which is pro-cyclic, and is
  described by the action of a topological generator of $I_x/P_x$. The
  group $I_x$ is the image of the Galois group of the field of fractions of
  the strict henselization of $\mathbb P^1_{\bar\F_q}$ at $x$; for a dense
  open $X'\subseteq X$ this map factors through $\pi_1^{\rm geom}(X')$, so the
  surjection of \ref{it:prelim:E1} maps the inertia and wild inertia groups of
  $X'$ at $x$ onto those of $X$, and $\mathcal F|_{X'}$ has at $x$ the local
  monodromy of $\mathcal F$; in particular it is tame at $x$ if $\mathcal F$
  is. If
  $\mathcal F$ is the restriction of a lisse sheaf on a dense open
  $X_0\supseteq X$ with $x\in X_0(\bar\F_q)$, then
  $\rho_{\mathcal F}(I_x)=1$; in particular $\mathcal F$ is tame at $x$. A
  tensor product of sheaves tame at $x$ is tame at $x$. An isomorphism
  $\mu:X'\to X$ that extends to an automorphism $\bar\mu$ of $\mathbb P^1$
  maps the inertia groups of $X'$ at $x$ onto those of $X$ at $\bar\mu(x)$,
  so $\mu^*\mathcal F$ is tame at $x$ iff $\mathcal F$ is tame at
  $\bar\mu(x)$, and the local monodromy of $\mu^*\mathcal F$ at $x$ is that
  of $\mathcal F$ at $\bar\mu(x)$.
\item\label{it:prelim:E4} \emph{Kummer sheaves}
  \cite[Sommes trig., 1.2--1.8, in particular (1.7.6)]{SGA4half}. Let $p$ be
  odd, $\chi=\chi_p$, and $g\in\F_p[X]$ nonconstant. There is a lisse sheaf
  $\mathcal L_{\chi(g)}$ of rank one on $\mathbb A^1\setminus Z(g)$
  ($Z(g)$ the set of roots of $g$) with
  $\mathrm{Tr}(\mathrm{Frob}_z\mid\mathcal L_{\chi(g)})=\chi(g(z))$ for
  every $z\in\F_p$ with $g(z)\ne0$. Its monodromy representation is the
  character of order dividing $2$ through which $\pi_1$ acts on the finite
  \'etale covering $\{w^2=g(z)\}$ (the geometric Frobenius at $z$ acts on the
  fibre $\{\pm w_0\}$, $w_0^2=g(z)$, by the inverse of $w_0\mapsto w_0^p$, fixing $w_0$ if $g(z)$ is a square in $\F_p$ and swapping $\pm w_0$ otherwise, i.e.\ by the
  sign $\chi(g(z))=g(z)^{(p-1)/2}$). Hence every Frobenius eigenvalue is
  $\pm1$, so $\mathcal L_{\chi(g)}$ is pure of weight $0$, and $P_x$, a
  pro-$p$ group with $p$ odd, acts trivially, so $\mathcal L_{\chi(g)}$ is
  tame at every point of $\mathbb P^1\setminus(\mathbb A^1\setminus Z(g))$.
  More precisely, let $x$ be such a point and $m_x$ the multiplicity of $x$
  as a root of $g$, with $m_\infty=\deg g$. The local monodromy of
  $\mathcal L_{\chi(g)}$ at $x$ is trivial if $m_x$ is even and is the
  unique character of order $2$ of $I_x$ (a topological generator of
  $I_x/P_x$ acts by $-1$) if $m_x$ is odd. Indeed, near a root $x$ write
  $g=(z-x)^{m_x}u$ with $u(x)\ne0$; $u$ has a square root over the strict
  henselization at $x$, because $p$ is odd, so the covering is locally
  $w'^2=(z-x)^{m_x}$: split if $m_x$ is even, the tame Kummer covering of
  degree $2$ if $m_x$ is odd. At $\infty$ the same holds in the coordinate
  $1/z$, in which $g$ is $(1/z)^{-\deg g}$ times a unit. In particular
  $\mathcal L_{\chi(g)}$ is geometrically nontrivial unless $g$ is a constant
  times a square in $\bar\F_p[X]$.
\end{enumerate}

\begin{theorem}[Grothendieck--Lefschetz trace formula; {\cite[Rapport, Thm.~3.2; Sommes trig., (1.1.1)]{SGA4half}}]\label{thm:prelim:lefschetz}
Let $X$ be a separated scheme of finite type over $\F_q$ and $\mathcal F$ a
constructible $\overline{\mathbb Q}_\ell$-sheaf on $X$. Then
\[
  \sum_{x\in X(\F_q)}\mathrm{Tr}(\mathrm{Frob}_x\mid\mathcal F_{\bar x})
  =\sum_i(-1)^i\,\mathrm{Tr}\big(\mathrm{Frob}_q\mid
  H^i_c(X\otimes\bar\F_q,\mathcal F)\big).
\]
\end{theorem}

\noindent (The Rapport states this for $\mathbb Q_\ell$-sheaves, and Sommes
trig.\ (1.1.1) for $E_\lambda$-sheaves, $E_\lambda$ a finite extension of
$\mathbb Q_\ell$; a constructible $\overline{\mathbb Q}_\ell$-sheaf comes from
an $E_\lambda$-sheaf for some such $E_\lambda$
\cite[(1.1.1), (1.1.6)]{DeligneWeilII}. The sheaves to which we apply it, in
the proofs of \cref{thm:dl:H2} and of the moving-root moments, are in fact
$\mathbb Q_\ell$-sheaves.)

\begin{theorem}[Cohomology of open curves; {\cite[(1.4.1)(a), (1.4.1.3), (1.4.2.3)]{DeligneWeilII}}, {\cite[Arcata, IV, Thm.~6.1]{SGA4half}}]\label{thm:prelim:curves}
Let $X$ be a smooth connected curve over an algebraically closed field that
is not complete, and $\mathcal F$ a lisse $\overline{\mathbb Q}_\ell$-sheaf on
$X$ with generic stalk $V$, a representation of $\pi_1(X)$. Then
$H^0_c(X,\mathcal F)=0$ and $H^2_c(X,\mathcal F)\cong V_{\pi_1(X)}(-1)$,
where $V_{\pi_1(X)}$ is the space of coinvariants. Moreover
$H^i_c(X,\mathcal F)=0$ for $i\ge3$ (by the vanishing theorem
\cite[Arcata, IV, Thm.~6.1]{SGA4half}, $R^if_!$ vanishes for $i$ larger than
twice the dimension of the fibres of $f$; take $f:X\to\mathrm{Spec}\,k$).
If $X=X_0\otimes\bar\F_q$ for a smooth geometrically connected curve $X_0$
over $\F_q$ and $\mathcal F$ comes from a lisse sheaf on $X_0$, this
isomorphism is compatible with the actions of $\mathrm{Gal}(\bar\F_q/\F_q)$
\cite[(1.4.2.3)]{DeligneWeilII}.
\end{theorem}

\begin{theorem}[Deligne: bounds on weights under $R^if_!$; {\cite[Thm.~3.3.1, Cor.~3.3.4 and (3.3.10)]{DeligneWeilII}}]\label{thm:prelim:weights}
Let $f:X\to Y$ be a morphism of schemes of finite type over $\F_q$ and
$\mathcal F$ a sheaf on $X$ that is $\jmath$-mixed of weights $\le w$. Then
$R^if_!\mathcal F$ is $\jmath$-mixed of weights $\le w+i$ for every $i$. In
particular ($Y=\mathrm{Spec}\,\F_q$) every eigenvalue $\alpha$ of
$\mathrm{Frob}_q$ on $H^i_c(X\otimes\bar\F_q,\mathcal F)$ satisfies
$|\jmath(\alpha)|\le q^{(w+i)/2}$.
\end{theorem}

\noindent (Theorem 3.3.1 and Corollary 3.3.4 are stated for sheaves that are
mixed in the sense of \cite[(1.2.2)]{DeligneWeilII}, i.e.\ for every
$\jmath$; (3.3.10) states that they hold, with the same proof, for
$\jmath$-mixed sheaves. The theorem is applied to sheaves that are pure for
every $\jmath$: of weight $1$ in the proof of \cref{thm:dl:H2} and of the
first moments \eqref{eq:advance:rank-first}, of weight $2$ in the second
moments \eqref{eq:advance:rank-second}; so either form suffices.)

\begin{theorem}[Grothendieck--Ogg--\v Safarevi\v c formula; {\cite[Thm.~1]{Raynaud}}, {\cite[Exp.~X, \S7]{SGA5}}]\label{thm:prelim:gos}
Let $k$ be an algebraically closed field of characteristic $p$,
$X\subseteq\mathbb P^1_k$ a dense open subscheme and $\mathcal F$ a lisse
$\overline{\mathbb Q}_\ell$-sheaf on $X$. Then
\[
  \sum_i(-1)^i\dim H^i_c(X,\mathcal F)
  =\mathrm{rk}(\mathcal F)\cdot\big(2-\#(\mathbb P^1(k)\setminus X(k))\big)
  -\sum_{x\in\mathbb P^1(k)\setminus X(k)}\mathrm{Swan}_x(\mathcal F).
\]
\end{theorem}

\noindent (Raynaud states the formula for constructible $\ell$-torsion
sheaves; the $\ell$-adic form above is the one quoted and used in
\cite[proof of Lemma~10.1.12]{KatzSarnak}. Here
$2-\#(\mathbb P^1(k)\setminus X(k))$ is the Euler characteristic of $X$ with
$\overline{\mathbb Q}_\ell$ coefficients. We use it only for sheaves that are
tame at every point of $\mathbb P^1(k)\setminus X(k)$, for which all Swan
terms vanish, \ref{it:prelim:E3}.)

\begin{theorem}[Katz--Sarnak: the family $y^2=h(x)(x-t)$ and its monodromy {\cite[\S10.1]{KatzSarnak}}]\label{thm:prelim:ks}
Let $k$ be a finite field of odd characteristic, $g\ge1$, and $h\in k[X]$ a
polynomial of degree $2g$ with $2g$ distinct roots in $\bar k$ (not
necessarily monic and not necessarily split: these are exactly the
hypotheses of \cite[10.1.1]{KatzSarnak}). Let
$U_h=\mathbb A^1_k\setminus Z(h)=\mathrm{Spec}\,k[T][1/h(T)]$, let
$\mathrm{pr}:\mathcal A_h\to U_h$ be the family of affine curves
$y^2=h(x)(x-t)$, $t\in U_h$, and
$\mathcal F_h=R^1\mathrm{pr}_!\overline{\mathbb Q}_\ell$ (Katz and Sarnak
construct a proper smooth model $\pi:\mathcal C_h\to U_h$ of this family whose
geometric fibres are connected curves of genus $g$, and show
$R^1\mathrm{pr}_!\overline{\mathbb Q}_\ell\cong
R^1\pi_*\overline{\mathbb Q}_\ell$, \cite[10.1.3--10.1.11]{KatzSarnak}).
Then:
\begin{enumerate}[label=\textup{(\roman*)}]
\item \cite[Lemma~10.1.9]{KatzSarnak} $\mathcal F_h$ is lisse of rank $2g$ on
  $U_h$ and pure of weight $1$;
\item \cite[Lemma~10.1.12]{KatzSarnak} $\mathcal F_h$ is tame at every point
  of $\mathbb P^1\setminus U_h$, the point $\infty$ included;
\item \cite[Lemma~10.1.15]{KatzSarnak} $\mathcal F_h$ is geometrically
  irreducible (indeed its geometric monodromy group is $\mathrm{Sp}(2g)$, see
  \textup{(v)});
\item (the identity in the proof of \cite[Lemma~10.1.15]{KatzSarnak}) for
  every finite extension $E$ of $k$ and every $t\in U_h(E)$,
  \[
    -\mathrm{Tr}(\mathrm{Frob}_{E,t}\mid\mathcal F_h)
    =\sum_{x\in E}\chi_E\big((x-t)h(x)\big),
  \]
  $\chi_E$ the quadratic character of $E^\times$ (extended by
  $\chi_E(0)=0$);
\item \cite[Lemma~10.1.9 and Theorem~10.1.16]{KatzSarnak} the cup product
  $\mathcal F_h\times\mathcal F_h\to R^2\mathrm{pr}_!\overline{\mathbb
  Q}_\ell|_{U_h}=\overline{\mathbb Q}_\ell(-1)$ is a perfect alternating
  pairing (a symplectic autoduality of $\mathcal F_h$ with values in
  $\overline{\mathbb Q}_\ell(-1)$), and the geometric monodromy group of
  $\mathcal F_h$, the Zariski closure of
  $\rho_{\mathcal F_h}(\pi_1^{\rm geom}(U_h))$, is the full symplectic group
  $\mathrm{Sp}(2g)$ of this pairing; in particular it is connected;
\item \cite[Lemma~10.1.13 and Remark~10.1.17]{KatzSarnak} the local
  monodromy \ref{it:prelim:E3} of $\mathcal F_h$ at each zero of $h$ is a
  unipotent pseudoreflection: a topological generator of the tame inertia
  acts by an element $u$ with $u-1$ of rank one and $(u-1)^2=0$; at
  $\infty$ it acts by $-u$ for such an element $u$. In particular the local
  monodromy of $\mathcal F_h$ at every point of $\mathbb P^1\setminus U_h$ is
  not scalar.
\end{enumerate}
\end{theorem}

\noindent The identity (iv) is also immediate from
\cref{thm:prelim:lefschetz,thm:prelim:curves}: the stalk of $\mathcal F_h$ at
a geometric point $\bar t$ over $t$ is
$H^1_c(\mathcal A_{h,\bar t},\overline{\mathbb Q}_\ell)$ (base change for
$R\mathrm{pr}_!$, \cite[Arcata, IV, Thm.~5.4]{SGA4half}); the fibre
$\mathcal A_{h,t}$ is a smooth geometrically irreducible affine curve
($(x-t)h(x)$ is squarefree of odd degree), so $H^0_c=0$ and
$H^2_c=\overline{\mathbb Q}_\ell(-1)$; and
$\#\mathcal A_{h,t}(E)=\sum_{x\in E}\big(1+\chi_E((x-t)h(x))\big)$. Katz and
Sarnak do not assume $h$ monic: whatever its leading coefficient, the curve
$y^2=h(x)(x-t)$ has a single point at infinity, at which their affine model $V^2=Uh^{\rm rev}(U)(1-TU)$, $h^{\rm rev}(U)=U^{2g}h(1/U)$, is smooth
\cite[10.1.1]{KatzSarnak}.

Items (v) and (vi) are what the moving-root second moments
\eqref{eq:advance:rank-second} need beyond (i)--(iv), after
\cref{lem:dl:sheaf}(v),(vi) has deduced the corresponding pairing and local-monodromy statements for polynomials with an odd number of roots by composing with inversion at a root: in the proof of \cref{lem:moving-root-moments} the
pairing of (v) identifies each moving-root sheaf with its dual, and the local
monodromy at a single root, from (vi) at a root of a moving-root sheaf and
from \ref{it:prelim:E4} at a root of the twist, shows that two of them are
isomorphic after a quadratic twist only if their root sets agree and the twist
is trivial (compare the invariant of $\mathcal F_h\otimes\mathcal F_h$
computed in the proof of \cite[Lemma~10.1.15]{KatzSarnak}). Only the pairing
in (v) and the monodromy at the zeros of $h$ in (vi) are used, and the
diagonal main term is computed elementarily. Katz and Sarnak deduce (vi), and
again (iii), from the description of $\mathcal F_h$ as a middle convolution (Katz,
\emph{Rigid local systems}, 3.3.6); their Lemma 10.1.13 proves directly that
the local monodromy at each zero of $h$ is trivial or a unipotent
pseudoreflection, and Remark 10.1.17 (printed p.~300) adds that it is
nontrivial and describes the monodromy at $\infty$. For $g=1$ and
$h(x)=x(x-1)$, $\mathcal F_h$ is the first cohomology of the Legendre family,
whose local monodromies at $t=0,1$ are transvections and at $t=\infty$ minus
a transvection.

\subsection{A modular coefficient}\label{subsec:external:modular}

The second moment of the neighbourhood count in the variance credits
(\cref{lem:record:weighted-moment}) contains the moment of one graph on five
labels, which a translation and a dilation reduce to the twisted second moment
$\sum_\lambda\chi(\lambda)T_\lambda^2$ of the traces $T_\lambda$ of the Legendre
family of elliptic curves. Weil's bound $|T_\lambda|\le2\sqrt p$ alone gives
$O(p^2)$ for this sum. Ahlgren and Ono evaluate it exactly, through a Gaussian
hypergeometric series and a cusp form of weight four, and the evaluation is
$O(p^{3/2})$.

\begin{theorem}[Ahlgren--Ono: a character-sum identity via the cusp form $\eta(2z)^4\eta(4z)^4$; Theorems~2 and~6, Lemma~2.2 {\cite{AhlgrenOno2000}}]\label{thm:ahlgren-ono}
Let $p$ be an odd prime, $\chi$ the quadratic character of $\Fp$ with
$\chi(0)=0$, $T_\lambda=\sum_{y\in\Fp}\chi\big(y(y-1)(y-\lambda)\big)$ for
$\lambda\in\Fp$, and let the integers $a(n)$ be the coefficients of
\[
 \sum_{n\ge1}a(n)q^n=q\prod_{n\ge1}(1-q^{2n})^4(1-q^{4n})^4=\eta(2z)^4\eta(4z)^4,
 \qquad q=e^{2\pi iz},
\]
where $\eta(z)=q^{1/24}\prod_{n\ge1}(1-q^n)$ is Dedekind's eta function. Then
\begin{enumerate}[label=\textup{(\alph*)}]
\item $\sum_{\lambda\in\Fp}\chi(\lambda)T_\lambda^2=-p-a(p)$;
\item $|a(p)|\le2p^{3/2}$.
\end{enumerate}
\end{theorem}

\noindent\emph{Source and normalization.} Ahlgren and Ono extend every
multiplicative character of $\Fp$, the trivial character $\varepsilon$ included,
by the value $0$ at $0$; they write $\phi$ for the quadratic character and
${}_{n+1}F_n(x)$ for Greene's Gaussian hypergeometric series whose upper
characters are all quadratic and whose lower characters are all trivial. Their
formula (2.1), from Greene, is
$p\,{}_2F_1(\lambda)=\varepsilon(\lambda)\phi(-1)\sum_x\phi\big(x(1-x)(1-\lambda x)\big)$.
For $\lambda\ne0$ the term $x=0$ vanishes, and the substitution $x=1/y$ turns
$x(1-x)(1-\lambda x)$ into $y^{-4}\,y(y-1)(y-\lambda)$; so
$p\,{}_2F_1(\lambda)=\chi(-1)T_\lambda$. The term $\lambda=0$ of (a) vanishes
because $\chi(0)=0$, so no identity is needed there. Their Lemma~2.2,
$p^3\,{}_4F_3(1)=p^2\sum_\lambda\phi(\lambda)\,{}_2F_1(\lambda)^2$, therefore reads
$p^3\,{}_4F_3(1)=\sum_\lambda\chi(\lambda)T_\lambda^2$, and their Theorem~6,
$p^3\,{}_4F_3(1)=-a(p)-p$ with the coefficients $a(n)$ of their (1.5), is (a). In
the proof of their Theorem~2 they observe that $\eta(2z)^4\eta(4z)^4$ is a cusp
form of weight four on $\Gamma_0(8)$ spanning a space of dimension one, hence a
normalized Hecke eigenform, and Deligne's theorem \cite{DeligneWeilI} gives (b);
their Theorem~2 is the consequence $|{}_4F_3(1)+p^{-2}|\le2p^{-3/2}$. Both
statements hold for every odd prime, with no congruence condition and no lower
bound on $p$; the applications take $p\equiv3\pmod4$. The theorem, lemma and
equation numbers here refer specifically to the authenticated
\href{https://uva.theopenscholar.com/files/ken-ono/files/049.pdf}{author manuscript},
not to an independently checked numbering of the published typesetting. In that manuscript (1.5) is on
printed page~3, (2.1) and Lemma~2.2 are on page~5, Theorem~6 is on page~7 and
the proof of Theorem~2 is on page~11; the
copy consulted (SHA-256 \texttt{e8bcbe78\ldots74b96a}) and the reduction of the
graph moment to (a) are recorded in
\path{research/roadmap_2026_10_18/item17/AHLGREN_ONO_SOURCE_REVIEW.md}. Finite
evaluations of both sides at small primes are normalization checks only; they
are not inputs of the proof.

\subsection{Sources of the standard inputs}\label{subsec:external:sources}

\emph{Weil's bound.} Part (a) of \cref{thm:weil} is the case of the prime
field and the quadratic character of \cite[Theorem~11.23]{IwaniecKowalski}
(with $n=1$ there): if $\chi$ is a multiplicative character of order $d>1$
of a finite field $F$ with $q$ elements and $f\in F[X]$ has $m$ distinct
roots and is not a constant times a $d$-th power over $\bar F$, then
$\big|\sum_{x\in F}\chi(f(x))\big|\le(m-1)\sqrt q$. Iwaniec and Kowalski
take it from \cite[Theorem~2C$'$]{Schmidt76}; the bound goes back to Weil
\cite{Weil48,Weil49}. A product of $t\ge1$ distinct linear factors has $t$
simple roots, so it is not a square. Part (b) of \cref{thm:weil} is deduced
from part (a) where it is stated and needs no further input.

\emph{Primes in the residue classes modulo four.} \Cref{thm:bmor} is
\cite[Theorem~1.2]{BMOR18} for the modulus $q=4$: for $q\ge3$ and $a$ coprime
to $q$, $|\theta(x;q,a)-x/\varphi(q)|<c_\theta(q)\,x/\log x$ for all
$x\ge x_\theta(q)$, and Table~1 (p.~432 of the published version) gives
$c_\theta(4)=0.0004822$ and $x_\theta(4)=4\,800\,162\,889$. The preprint
\texttt{arXiv:\allowbreak1802.00085v2} lists different values for $q=4$; only
the published ones are used.

\emph{Kassabov's inequality.} \Cref{thm:kassabov} is
\cite[Lemma~4.2(b)]{Kassabov} (printed p.~15 of
\texttt{arXiv:\allowbreak0911.1983}); the companion bound for
$\cos^2\angle(V_1\cap V_2,V_3)$, Lemma 4.1(b) there, is not used. Kassabov's
angle between two closed subspaces is the Friedrichs angle of
\cref{def:prelim:angle} \cite[Definition~3.2]{Kassabov}, defined when neither
subspace contains the other. He states the lemma for numbers
$\eps_{ij}=\cos\alpha_{ij}$ with $0<\alpha_{ij}\le\angle(V_i,V_j)$; in finite
dimension such angles are positive \cite[Corollary~3.7]{Kassabov}, and
$\alpha_{ij}=\angle(V_i,V_j)$ gives the form of \cref{thm:kassabov}. If one of
the three subspaces contains another, \cref{def:prelim:angle} gives that
pair the cosine $0$, and the inequality holds directly. If $V_1\subseteq V_3$,
the left side is $\cos\angle(V_1,V_2\cap V_3)$, and the vectors admitted in its
supremum are admitted for $\cos\angle(V_1,V_2)$, both being orthogonal to
$V_1\cap V_2\cap V_3=V_1\cap V_2$; so it is at most $\eps_{12}$, which is at
most the right side since $\eps_{13}=0$. The case $V_2\subseteq V_3$ is
symmetric, and in the remaining cases one of $V_1\cap V_3$, $V_2\cap V_3$
contains the other, so the left side is $0$.

\paragraph{Consequences and what to read next.}
This appendix assumes nothing from the rest of the paper, except \cref{def:prelim:angle} for the
matching of Kassabov's statement. The remarks after \cref{thm:prelim:gl2}
and \cref{thm:prelim:ks} use only the cited sources and
\cref{thm:prelim:lefschetz,thm:prelim:curves}.
